%% L'escalier des préfixes du Système international, du nano au giga.
%% Les puissances de 10 sont la réponse : elles passent par \rappel (variante muette dans les énoncés).
\documentclass[tikz,border=4pt]{standalone}
\usepackage{liaisons}
\begin{document}
\begin{tikzpicture}[x=1cm,y=1cm,
  marche/.style={draw=solideB, line width=1pt, fill=solideB!8, rounded corners=1.5pt,
                 minimum width=1.32cm, minimum height=0.62cm, inner sep=1pt, font=\small\bfseries},
  base/.style={draw=solideA, line width=1.2pt, fill=solideA!12, rounded corners=1.5pt,
               minimum width=1.32cm, minimum height=0.62cm, inner sep=1pt, font=\small\bfseries},
  exp/.style={font=\footnotesize, text=black!70, anchor=north, inner sep=2pt}]

\foreach \i/\sym/\nom/\pow in {%
  0/n/nano/{10^{-9}}, 1/$\mu$/micro/{10^{-6}}, 2/m/milli/{10^{-3}},
  3/{\footnotesize unité}/{}/{10^{0}}, 4/k/kilo/{10^{3}}, 5/M/méga/{10^{6}}, 6/G/giga/{10^{9}}} {
  \pgfmathsetmacro\x{1.42*\i}
  \ifnum\i=3
    \node[base] (n\i) at (\x,0) {\sym};
  \else
    \node[marche] (n\i) at (\x,0) {\sym};
  \fi
  \node[font=\scriptsize, text=black!65, anchor=south, inner sep=3pt] at (\x,0.31) {\nom};
  \node[exp] at (\x,-0.33) {\rappel{$\pow$}};
}

%% flèches de passage, avec le facteur
\foreach \a/\b in {0/1,1/2,2/3,3/4,4/5,5/6} {
  \draw[-{Stealth[length=5pt,width=4pt]}, line width=0.9pt, draw=black!45]
     ([yshift=-1.05cm]n\a.south) -- ([yshift=-1.05cm]n\b.south);
}
\foreach \i in {0,...,5} {
  \pgfmathsetmacro\x{1.42*\i+0.71}
  \node[font=\scriptsize, text=black!60, anchor=north, inner sep=1.5pt] at (\x,-1.42) {$\times10^{3}$};
}
\node[font=\scriptsize\itshape, text=black!55, anchor=north, inner sep=3pt] at (4.26,-1.88)
  {vers la droite on multiplie par mille, vers la gauche on divise par mille};
\end{tikzpicture}
\end{document}
